\section*{الفصل \ref{ch:b1:potential-capacitors} --- الكمون الكهربائي؛ النواقل والمكثفات}

\begin{solution}{exo:b1:potential-capacitors:1}
$V = 8.99 \times 10^9 \times 10^{-6}/1 = \qty{9.0}{kV}$؛ و $W = qV = \qty{9.0}{mJ}$؛
وإذا أُطلقت: $\tfrac12mv^2 = \qty{9}{mJ}$، ومنه $v = \sqrt{18} = \qty{4.2}{m/s}$.
\end{solution}

\begin{solution}{exo:b1:potential-capacitors:2}
$E = 100/10^{-3} = \qty{1e5}{V/m}$؛ و $E_k = eU = \qty{100}{eV} = \qty{1.6e-17}{J}$،
ومنه $v = \sqrt{2E_k/m_e} = \qty{5.9e6}{m/s}$؛ وبتسارع منتظم يكون
$t = 2d/v = \qty{3.4e-10}{s}$.
\end{solution}

\begin{solution}{exo:b1:potential-capacitors:3}
$E_A = 10/0.002 = \qty{5}{kV/m}$ و $E_B = 10/0.008 = \qty{1.25}{kV/m}$. ويتجه
المجال نحو الكمون المتناقص؛ ويتسارع إلكترون ($q < 0$)
في الجهة الأخرى، نحو $V$ المتزايد.
\end{solution}

\begin{solution}{exo:b1:potential-capacitors:4}
$C = 4\pi\varepsilon_0R = 1.11 \times 10^{-10} \times 0.15 = \qty{17}{pF}$؛ و $Q = CV =
\qty{3.3}{\micro C}$؛ و $E = V/R = 2 \times 10^5/0.15 = \qty{1.3}{MV/m}$. ويقع الانهيار
عند $V = E_{\max}R = 3 \times 10^6 \times 0.15 = \qty{450}{kV}$ --- ولذا تستعمل الآلات
الكبيرة كرات كبيرة.
\end{solution}

\begin{solution}{exo:b1:potential-capacitors:5}
كل عنصر من الحلقة على البعد $\sqrt{R^2 + z^2}$: ومنه $V =
Q/4\pi\varepsilon_0\sqrt{R^2 + z^2}$ و $E_z = -\dd V/\dd z = Qz/4\pi\varepsilon_0(R^2 +
z^2)^{3/2}$. وفي القرص: $\dd q = 2\pi\sigma a\,\dd a$ على البعد $\sqrt{a^2 + z^2}$، ومنه $V =
\dfrac{\sigma}{2\varepsilon_0}\displaystyle\int_0^R\dfrac{a\,\dd a}{\sqrt{a^2 + z^2}} =
\dfrac{\sigma}{2\varepsilon_0}\bigl(\sqrt{z^2 + R^2} - z\bigr)$، ويكون $-\dd V/\dd z =
\dfrac{\sigma}{2\varepsilon_0}\Bigl(1 - \dfrac{z}{\sqrt{z^2 + R^2}}\Bigr)$.
\end{solution}

\begin{solution}{exo:b1:potential-capacitors:6}
$V(r) = V(R) + \displaystyle\int_r^R\dfrac{Qr'\,\dd r'}{4\pi\varepsilon_0R^3} =
\dfrac{Q}{4\pi\varepsilon_0R} + \dfrac{Q(R^2 - r^2)}{8\pi\varepsilon_0R^3} = \dfrac{Q}{4\pi\varepsilon_0}
\dfrac{3R^2 - r^2}{2R^3}$. وفي الذهب: $V(0) = 1.5 \times 8.99 \times 10^9 \times 1.26 \times
10^{-17}/7 \times 10^{-15} = \qty{24}{MV}$: فيحتاج بروتون إلى \qty{24}{MeV}، وجسيم ألفا
($2e$) إلى \qty{48}{MeV} ليبلغ المركز وإلى \qty{32}{MeV} ليلمس
السطح. وترتدّ جسيمات ألفا ذات \qty{5}{MeV} عند رذرفورد عند $r$ حيث $2eV =
\qty{5}{MeV}$: أي $V = \qty{2.5}{MV}$ و $r = \qty{45}{fm}$ --- وهو خارج النواة
بكثير، ولهذا نجح تحليله بالشحنة النقطية.
\end{solution}

\begin{solution}{exo:b1:potential-capacitors:7}
على المحور: $2p/4\pi\varepsilon_0r^3 = 2 \times 8.99 \times 10^9 \times 6.2 \times 10^{-30}/10^{-27} =
\qty{1.1e8}{V/m}$؛ وعموديًّا عليه النصف: \qty{5.6e7}{V/m}. والعزم المزدوج $pE =
\qty{6.2e-24}{N.m}$؛ وفرق الطاقة $2pE = \qty{1.2e-23}{J}$ مقابل $k_BT =
\qty{4.1e-21}{J}$: أي أصغر $330$ مرة --- فلا مطابقة صافية إلا بنسبة $0.1\%$.
وأيون الصوديوم على \qty{0.3}{nm}: $E = e/4\pi\varepsilon_0r^2 = \qty{1.6e10}{V/m}$ و $-pE
= \qty{-1.0e-19}{J} = \qty{-0.6}{eV}$، أي $24\,k_BT$: فتتشبث جزيئات الماء
بالأيون --- وهي الإماهة.
\end{solution}

\begin{solution}{exo:b1:potential-capacitors:8}
$Q = 4\pi\varepsilon_0R^2E = 10^{-4} \times 3 \times 10^6/8.99 \times 10^9 = \qty{33}{nC}$
و $V = ER = \qty{30}{kV}$. وفي الكرتين الموصولتين: $E_1 = V/R_1 = \qty{0.2}{MV/m}$
و $E_2 = V/R_2 = \qty{20}{MV/m} > \qty{3}{MV/m}$: فتؤيّن الكرة الصغيرة
الهواء وتنزف الشحنة --- وهو \omterm{ex:b1:potential-capacitors:point}{أثر الأطراف}.
\end{solution}

\begin{solution}{exo:b1:potential-capacitors:9}
$C = \varepsilon_r\varepsilon_0S/d = 2.2 \times 8.85 \times 10^{-12}/10^{-5} = \qty{2.0}{\micro F}$؛
واللفّ يغيّر الشكل لا المساحة المتقابلة: فتبقى $C$ هي هي (ومع شريحة
ثانية بحيث تقابل كل رقيقة الأخرى من الوجهين، تتضاعف).
و $W = \tfrac12 \times 2 \times 10^{-6} \times 400^2 = \qty{0.16}{J}$. وفي مزيل الرجفان:
$Q = 150 \times 10^{-6} \times 2000 = \qty{0.30}{C}$ و $W = \tfrac12CU^2 = \qty{300}{J}$
و $P = 300/0.005 = \qty{60}{kW}$.
\end{solution}

\begin{solution}{exo:b1:potential-capacitors:10}
(أ) $E = Q/4\pi\varepsilon_0r^2$ و $U = \dfrac{Q}{4\pi\varepsilon_0}\Bigl(\dfrac1{R_1} -
\dfrac1{R_2}\Bigr)$ و $C = 4\pi\varepsilon_0R_1R_2/(R_2 - R_1)$؛ ومع $R_2 - R_1 = d \ll
R_1$ يكون $R_1R_2 \approx R^2$ و $C \approx 4\pi R^2\varepsilon_0/d = \varepsilon_0S/d$؛ ومع $R_2 \to
\infty$ يكون $C \to 4\pi\varepsilon_0R_1$. (ب) $C/L = 2\pi\varepsilon_r\varepsilon_0/\ln(b/a) = 2\pi \times 2.3
\times 8.85 \times 10^{-12}/\ln5 = \qty{79}{pF/m}$، وهي الرتبة الصحيحة. (ج) $E(a) =
U/[a\ln(b/a)] = 1000/(5 \times 10^{-4} \times 1.61) = \qty{1.2}{MV/m}$: أي دون
\qty{20}{MV/m} للبولي إيثيلين، فهو آمن؛ أما في الهواء فكان على الحدّ.
\end{solution}

\begin{solution}{exo:b1:potential-capacitors:11}
(أ) $W = Q^2/2C = Q^2/8\pi\varepsilon_0R$. (ب) و $R = e^2/8\pi\varepsilon_0m_ec^2 = 2.31 \times
10^{-28}/(2 \times 8.19 \times 10^{-14}) = \qty{1.4e-15}{m}$ (وأما ``نصف القطر
الكلاسيكي'' المعتاد $e^2/4\pi\varepsilon_0m_ec^2 = \qty{2.8}{fm}$ فيسقط العامل $\tfrac12$). (ج) و $W =
Q^2/2C = Q^2d/2\varepsilon_0S$ يتضاعف: ومنه $\Delta W = Q^2d/2\varepsilon_0S = F\,d$ مع $F =
Q^2/2\varepsilon_0S = \sigma^2S/2\varepsilon_0$، كما وُجد مباشرةً. (د) وعند $U$ ثابت يكون $W =
\varepsilon_0SU^2/2d$ فينتصف: $\Delta W = -W/2$؛ ويؤدي الساحب مع ذلك الشغل $W/2$
(بمكاملة $F = \varepsilon_0SU^2/2x^2$ من $d$ إلى $2d$)؛ والنصفان معًا، أي $W$
كله، يعودان إلى البطارية حين تعود إليها الشحنة $CU/2$ عند
التوتر $U$.
\end{solution}

\begin{solution}{exo:b1:potential-capacitors:12}
(أ) الكرة التي نصف قطرها $r$ تحمل $q = Qr^3/R^3$ عند كمون سطحي
$q/4\pi\varepsilon_0r$؛ وإضافة القشرة $\dd q = 3Qr^2\,\dd r/R^3$ تكلّف $\dd W =
q\,\dd q/4\pi\varepsilon_0r = 3Q^2r^4\,\dd r/4\pi\varepsilon_0R^6$؛ وبالمكاملة: $W =
3Q^2/20\pi\varepsilon_0R$. (ب) و $Q = 92e = \qty{1.47e-17}{C}$: ومنه $W = 3 \times 2.17 \times
10^{-34}/(20\pi \times 8.85 \times 10^{-12} \times 7.4 \times 10^{-15}) = \qty{1.6e-10}{J}
= \qty{990}{MeV}$. (ج) ولكل شظية: الشحنة $Q/2$ ونصف القطر $R/2^{1/3}$،
والطاقة $\tfrac14 2^{1/3}W$؛ ولشظيتين: $2^{1/3}W/2 = 0.63W$: فالمحرَّر
$0.37 \times 990 = \qty{370}{MeV}$. (د) والذي يعارضه هو ``التوتر السطحي'' النووي:
فللشظيتين سطح أكبر من سطح نواة واحدة، وذلك يكلّف نحو
\qty{170}{MeV}؛ ويظهر الباقي، وهو قرب \qty{200}{MeV}، في معظمه
طاقةً حركية للشظيتين وهما تتطايران تحت تنافرهما
المتبادل.
\end{solution}

\begin{solution}{pb:b1:potential-capacitors:1}
\textbf{1.} الثقل $\tfrac43\pi r^3\rho g$ نحو الأسفل، والجرّ $6\pi\eta rv$ نحو الأعلى،
والدافعة $\tfrac43\pi r^3\rho_{\text{الهواء}}g$ نحو الأعلى --- وهي $1.2/900 = 0.13\%$ من
الثقل، فتُهمَل. وعند \omterm{prop:b1:newton-dynamics:terminal}{السرعة الحدية}: $\tfrac43\pi r^3\rho g = 6\pi\eta rv_1$، ومنه $v_1 =
2\rho gr^2/9\eta$.

\textbf{2.} $v_1 = \qty{1.0e-4}{m/s}$: ومنه $r = \sqrt{9\eta v_1/2\rho g} = \sqrt{9 \times
1.8 \times 10^{-5} \times 10^{-4}/(2 \times 900 \times 9.8)} = \qty{9.6e-7}{m}$، أي
ميكرومتر؛ و $m = \tfrac43\pi r^3\rho = \qty{3.3e-15}{kg}$.

\textbf{3.} $\tau = m/6\pi\eta r = 3.3 \times 10^{-15}/3.3 \times 10^{-10} = \qty{1e-5}{s}$؛
والمسافة $\sim v_1\tau = \qty{1}{nm}$. فلا تُرى أبدًا.

\textbf{4.} $\mathrm{Re} = 1.2 \times 10^{-4} \times 9.6 \times 10^{-7}/1.8 \times 10^{-5} =
6 \times 10^{-6} \ll 1$: فقانون ستوكس يصحّ.

\textbf{5.} $eE = 1.6 \times 10^{-19} \times 3 \times 10^5 = \qty{4.8e-14}{N}$ مقابل
$mg = \qty{3.2e-14}{N}$: فهما متقاربان في قطرات ميكرومترية وحدها. وقطرة قطرها \qty{10}{\micro m}
تزن أكثر بألف مرة؛ وكانت \omterm{def:b1:electrostatics-gauss:charge}{شحنة أولية} واحدة لتغيّر
سرعتها بنسبة $0.1\%$، وهي غير مرئية.

\textbf{6.} إنها الحركة البراونية: فقصف جزيئات الهواء يزعزع
القطرة عشوائيًّا؛ فتتشتت التوقيتات، مما يحدّ من الدقة
ويضع حدًّا أدنى لحجم القطرة المفيد.

\textbf{7.} $E = U/d = 5000/0.016 = \qty{3.1e5}{V/m}$؛ والصفيحتان أعرض بكثير من
$d$، فيعطي تراكب المستويين مجالًا منتظمًا؛ والأسطح
المتساوية الكمون مستويات أفقية، $V = Uz/d$.

\textbf{8.} $C = \varepsilon_0\pi(0.1)^2/0.016 = \qty{17}{pF}$؛ و $Q = CU = \qty{87}{nC}$؛
و $W = \tfrac12CU^2 = \qty{0.22}{mJ}$.

\textbf{9.} الصفيحة العليا موجبة: فيتجه $\vect E$ نحو الأسفل؛ وعلى القوة أن
تتجه نحو الأعلى، ومنه $q < 0$. و $|q| = mgd/U_0 = 3.3 \times 10^{-15} \times 9.8 \times 0.016/
1080 = \qty{4.8e-19}{C} = 3e$: ومنه $q = -3e$.

\textbf{10.} في الصعود بسرعة ثابتة: $|q|E = mg + 6\pi\eta rv_2$ و $mg =
6\pi\eta rv_1$، ومنه $|q|E = 6\pi\eta r(v_1 + v_2)$ و $|q| = 6\pi\eta r(v_1 +
v_2)d/U$. و $v_2 = (|q|E - mg)/6\pi\eta r = (1.5 \times 10^{-13} - 3.2 \times 10^{-14})/
3.3 \times 10^{-10} = \qty{3.6e-4}{m/s}$: أي \qty{1}{mm} في \qty{2.8}{s}.

\textbf{11.} التوازن حكم على انعدام حركة ذرّة سابحة مرتجفة
قد تتغير شحنتها في أثناء ذلك؛ أما التوقيتات على مسافة ثابتة فيمكن
تكرارها كما نشاء، ويثبت $r$ مرة واحدة إلى الأبد بواسطة $v_1$، فتعطي قطرة واحدة
عندئذ سلسلة كاملة من الشحن.

\textbf{12.} $|q|U = 3 \times 1.6 \times 10^{-19} \times 5000 = \qty{2.4e-15}{J}$
(أي \qty{15}{keV}) عبر الفجوة؛ والجرّ على امتداد \qty{1}{cm}: $(|q|E - mg) \times 0.01 =
\qty{1.2e-15}{J}$، تتحول حرارةً في الهواء.

\textbf{13.} بالقسمة على $1.60$: $3.00$ ثم $5.01$ ثم $4.01$ ثم $1.99$ ثم $6.02$ ---
أي الأعداد الصحيحة $3, 5, 4, 2, 6$.

\textbf{14.} النسب $q/n$: $1.600, 1.602, 1.605, 1.595, 1.605$؛ والمتوسط $1.601$،
والانحراف المعياري $0.004$، وارتياب المتوسط $0.004/\sqrt5 =
0.002$: ومنه $e = (1.601 \pm 0.002) \times 10^{-19}$~كولوم.

\textbf{15.} تؤيّن الأشعة السينية الهواء؛ فتلتقط القطرة الأيونات واحدًا
تلو آخر، ويحمل كلٌّ منها عددًا صحيحًا من $e$.

\textbf{16.} $r \propto \eta^{1/2}$ انطلاقًا من $v_1$، ومنه $m \propto r^3 \propto \eta^{3/2}$، و
$|q| = mgd/U_0 \propto \eta^{3/2}$: فخطأ $1\%$ في $\eta$ يعطي $1.5\%$ في $e$. وكانت
$\eta$ عند ميليكان أدنى بنسبة $0.4\%$، فكانت $e$ عنده أدنى بنسبة $0.6\%$.

\textbf{17.} لا تأتي الكواركات وحدها أبدًا: فهي محبوسة داخل الهادرونات،
ويحمل كل جسم حرّ شحنة صحيحة.

\textbf{18.} $N_A = F/e = 96485/1.602 \times 10^{-19} = \qty{6.02e23}{mol^{-1}}$.

\textbf{19.} $m_e = 1.602 \times 10^{-19}/1.76 \times 10^{11} = \qty{9.1e-31}{kg}$: أي أخفّ $1840$
مرة من ذرة الهيدروجين.

\textbf{20.} مستويات أفقية بين الصفيحتين، تنتفخ حول
الحافتين؛ وخطوط مجال شاقولية تتفرّع نحو الخارج عند الحواف على عرض
من رتبة $d = \qty{1.6}{cm}$ --- وهو صغير أمام القطر \qty{20}{cm}، و
القطرة تُرقَب في المركز.

\textbf{21.} $E_p = qUz/d + mgz$، وهي خطية في $z$: فلا أصغري؛ وتقتضي
قطرة ساكنة انعدام الميل، أي $q = -mgd/U$، ويكون التوازن
عندئذ لامباليًا --- فيدوم أيّ انسياق متبقٍّ، ولهذا يصعب الحكم على
التوازن.

\textbf{22.} $2.4 \times 10^{-15}$ مقابل $2.2 \times 10^{-4}$~جول: أي $10^{-11}$. فالقطرة
لا تضطرب المجال.

\textbf{23.} المجال نحو الأعلى، والقوة على القطرة السالبة نحو الأسفل: فتسقط بالسرعة $(mg
+ |q|E)/6\pi\eta r = (3.2 \times 10^{-14} + 1.5 \times 10^{-13})/3.3 \times 10^{-10} =
\qty{5.6e-4}{m/s}$.

\textbf{24.} $F = 10^4 \times 1.6 \times 10^{-19} \times 3.1 \times 10^5 = \qty{5e-10}{N}$
مقابل $\qty{9.8e-9}{N}$: أي $5\%$ من الثقل --- فلا يمكن رفعها، ومع
$10^4$ شحنة يكون تغيّر شحنة واحدة $e$ غير مرئي. ورذاذ الزيت يعطي قطرات
صغيرة وكروية ولا تتبخر.

\textbf{25.} المقيس: $v_1$ و $v_2$ و $U$ (مع معلومية $d$ و $\eta$ و $\rho$)؛
والمستنتج $r$ ثم $m$ ثم $q$، التي وُجدت دائمًا مضاعفًا صحيحًا للمقدار
$e = (1.601 \pm 0.002) \times 10^{-19}$~كولوم (بنسبة $0.1\%$)؛ ومنها $N_A = F/e$ و
$m_e = e/(e/m_e)$.
\end{solution}
